\chapter{队伍优化方法}\label{chap:method}

本章详细介绍本文提出的宝可梦队伍优化方法。首先阐述融合双属性组合与特性修正的属性覆盖面优化算法，然后介绍基于贝叶斯优化的速度线阈值预测模型，最后给出面向队伍整体协同的多目标遗传算法框架。

\section{属性覆盖面优化算法}

\subsection{问题建模}

设队伍 $T = \{P_1, P_2, \ldots, P_6\}$ 包含六只宝可梦，每只宝可梦 $P_i$ 具有技能集合 $M_i = \{m_{i1}, m_{i2}, m_{i3}, m_{i4}\}$。当使用技能 $m$ 攻击防御宝可梦 $D$ 时，伤害倍率为：

\begin{equation}\label{eq:damage-mult}
    f(m, D) = \prod_{t \in \text{types}(D)} \text{eff}(\text{type}(m), t) \times \text{stab}(m, P_i) \times \text{ability\_mod}(P_i, D)
\end{equation}

其中 $\text{eff}(a, b)$ 为属性 $a$ 对属性 $b$ 的基础克制倍率，$\text{stab}(m, P_i)$ 为本系加成因子，$\text{ability\_mod}(P_i, D)$ 为特性修正项。

队伍的属性覆盖度 $C(T)$ 定义为：对目标宝可梦集合 $\mathcal{D}$，存在至少一个技能使其伤害倍率超过阈值 $\tau$（通常取 $\tau=1.0$，即能造成正常伤害即可）的宝可梦数量占比：

\begin{equation}\label{eq:coverage}
    C(T) = \frac{|\{D \in \mathcal{D} : \exists P_i \in T, \exists m \in M_i, f(m, D) \geq \tau\}|}{|\mathcal{D}|}
\end{equation}

\subsection{算法设计}

本文提出一种贪心-局部搜索混合算法（Greedy-Local Search Hybrid, GLSH）来最大化队伍的属性覆盖度。算法分为两个阶段：

\textbf{阶段一：贪心构建}。从空队伍开始，每次选取使当前覆盖度增量最大的宝可梦加入队伍，直到队伍满6只或覆盖度不再提升。

\textbf{阶段二：局部搜索优化}。对贪心构建的结果进行邻域搜索，每次尝试替换队伍中的一只宝可梦，若覆盖度提升则接受替换，直到达到局部最优。

算法伪代码如算法~\ref{alg:glsh}所示。

\begin{algorithm}[!htbp]
    \small
    \caption{GLSH 属性覆盖优化算法}\label{alg:glsh}
    \begin{algorithmic}[1]
        \State 输入：候选宝可梦池 $\mathcal{P}$，目标防御宝可梦集合 $\mathcal{D}$
        \State 输出：最优队伍 $T^*$
        \State 初始化空队伍 $T \gets \emptyset$
        \Repeat
            \State $P^* \gets \argmax_{P \in \mathcal{P} \setminus T} [C(T \cup \{P\}) - C(T)]$
            \If {$C(T \cup \{P^*\}) > C(T) \land |T| < 6$}
                \State $T \gets T \cup \{P^*\}$
            \EndIf
        \Until {$|T| = 6$ 或覆盖度不再提升}
        \Repeat
            \State 选择 $P_{out} \in T$, $P_{in} \in \mathcal{P} \setminus T$
            \State $T' \gets (T \setminus \{P_{out}\}) \cup \{P_{in}\}$
            \If {$C(T') > C(T)$}
                \State $T \gets T'$
            \EndIf
        \Until 达到局部最优
        \State \Return $T$
    \end{algorithmic}
\end{algorithm}

\section{速度线阈值预测模型}

\subsection{基于贝叶斯优化的速度线分析}

速度线的确定具有高度的赛制依赖性。本文提出基于贝叶斯优化（Bayesian Optimization）的速度线阈值预测方法。设速度线阈值 $\theta^*$ 为待优化参数，目标函数 $g(\theta)$ 为在给定阈值下队伍的综合表现评分：

\begin{equation}\label{eq:bayes-opt}
    \theta^* = \argmax_{\theta \in [\theta_{\min}, \theta_{\max}]} g(\theta)
\end{equation}

贝叶斯优化通过高斯过程（Gaussian Process）代理模型对 $g(\theta)$ 进行建模，利用期望提升（Expected Improvement, EI）作为采集函数，在探索（Exploration）与利用（Exploitation）之间取得平衡。

\subsection{关键速度门槛的提取}

通过对VGC 2022--2025赛季数据进行分析，本文提取了如表~\ref{tab:speed-tiers}所示的关键速度门槛。

\begin{table}[!htbp]
    \bicaption{VGC 2022--2025赛季关键速度门槛}{Key speed tier thresholds in VGC 2022--2025 seasons}
    \label{tab:speed-tiers}
    \centering
    \footnotesize
    \begin{tabular}{ccccl}
        \hline
        速度门槛 & 实际速度值 & 要求种族值 & 努力值投入 & 代表性宝可梦 \\
        \hline
        极速线 & $\geq$200 & $\geq$120 & 252 EV & 雷吉艾勒奇 \\
        高速线 & 160--199 & 90--119 & 252 EV & 振翼发 \\
        中速线 & 120--159 & 60--89 & 适量EV & 土地云 \\
        低速线 & <120 & <60 & 最低EV & 青铜钟 \\
        空间线 & <50 & <30 & 0 EV & 垒磊石 \\
        \hline
    \end{tabular}
\end{table}

\section{多目标遗传算法框架}

\subsection{编码方案与适应度函数}

本文采用NSGA-II（Nondominated Sorting Genetic Algorithm II）作为多目标优化框架。每一条染色体编码一个队伍的六只宝可梦及各自的道具、技能和努力值分配方案。

适应度评估包含三个子目标：
\begin{enumerate}
    \item $f_1$：属性覆盖度 $C(T)$（最大化）；
    \item $f_2$：速度配合度 $S(T)$，衡量队伍的速度线阶梯分布合理性（最大化）；
    \item $f_3$：攻防平衡度 $B(T)$，衡量队伍的物理/特殊输出和防御均衡性（最大化）。
\end{enumerate}

\subsection{非支配排序与拥挤度计算}

NSGA-II通过非支配排序将种群划分为多个前沿层级，并在同一前沿内通过拥挤度距离（Crowding Distance）保持解的多样性。最终输出帕累托最优解集，训练家可根据个人偏好从中选择最终队伍配置。

\section{本章小结}

本章从属性覆盖优化、速度线预测和队伍协同三个方面构建了宝可梦PVP队伍优化框架。GLSH算法有效提升了队伍的属性覆盖面，贝叶斯优化模型实现了对速度线阈值的精确预测，NSGA-II框架则提供了兼顾多个目标的帕累托最优队伍配置方案。
